
This chapter discusses how questions can be imperfect and how both humans and computers deal with those imperfect questions.



<<<<<<< Updated upstream

But instead, I want to focus on linguistic phenomena that \emph{causes} a question to \emph{become} unanswerable.  
%
This line of research began with what has come to be known as as ``the lightbulb paper'' from the title \question{Which linguist invented the lightbulb}, 
%
which highlights the problem: there’s a well-defined answer to the question \question{who invented the first commercially successful light bulb}, 
%
but it sure as heck wasn’t a linguist\dots answering that question would perhaps not be useful to a user.

There are two takeaways here: First that many early \abr{qa} systems accepted questions like \question{When did Pablo Picasso paint the Mona Lisa} and answer 1503, hundreds of years before Picasso was born.\footnote{Because he didn't paint it.}
%
Second, that many of the unanswerable questions in Natural Questions are unanswerable because of these phenomena.  
%
And it’s not just Natural Questions: Yu et al~\citep{} showed that this also shows up on Reddit.  It’s an important problem!

But we don’t actually know \emph{why} a particular natural question is unanswerable.  
%
Is it because Wikipedia lacks the information, because there’s a false assumption, because there’s actually dispute about the canonical answer to a question, or because it’s a matter of taste / opinion?  
%
Because Natural Questions was focused on factual \abr{qa} (a reasonable position at the time), we don’t have a single comprehensive dataset to address all of these issues (the numbers here are just for a small sample of \abr{nq}).



Yogesh Raut’s recreational thinking podcast is an example of this interrogation process: they spend several minutes interrogating individual questions, figuring out possibilities, and then justifying the answer.
%
So for example, if the question comes in: \underline{What linguist invented the lightbulb}, you might have an interaction that goes like this:
%
``Um, actually, most of the people credited with the development of the lightbulb were chemists, physicists, or engineers''
%
``Okay, smartbutt, who made the first lightbulb''

“Well, do you mean created the first electric light (Humphry Davy), the first filament inside a vacuum tube (Jospeh Swan), the first metal filament in a vacuum tube (Warren de la Rue), or the first commercially successful lightbulb with those properties (Thomas Edison)?”

“Hmm … I guess I meant who got the credit for it”

“Then that would be Thomas Edison”
=======
\section{But That Assumes a Good Question}
\label{sec:methods:presuppositions}

Even a fluent generated answer can be unhelpful when the question is problematic.  For example, some questions that tripped up early \abr{qa} systems were ``When did Leonardo paint the Sistine Chapel'' or ``Is the king of France bald''.  These questions contain what are called ``false presuppositions'': Michelangelo painted the Sistine Chapel, and France no longer has a king.  A model can treat these as familiar-looking questions, accept their false presuppositions, and blithely answer.

You might think that these questions are unrealistic, but questions with hidden assumptions are common in the medical domain.
For example, \citet{srikanth-24-pregnant} examine the question ``Can I breastfeed my baby when I am ill?''  A model can go along with a question's assumptions---like an improv actor saying ``yes and''~\citep{cho-20-spolin}---but a helpful answer may need to challenge them.  In the paper's example, the fuller answer addresses the belief that all illnesses pass through breast milk and explains that antibodies in breast milk can help protect the baby.

But how do you push back?  This requires coming up with a strategy.  Sometimes you're missing information.  If the question is ``what's the capital of Georgia'', you might need to ask which Georgia they mean.  You might need to ask for context: ``how old is the baby'', ``what were you sick with''.  These are called discourse moves, and they direct the conversation in different directions.  \citet{srikanth-26-discotrace} show how different communities favor different answering strategies, including clarifying a question or challenging its assumptions.

Given this, how do you get to a good answer?
>>>>>>> Stashed changes
